%%%PAGE 5 rot=0 legibility=fair summary=选5：热辐射与黑体辐射、电磁波谱、普朗克能量子与电子伏特、光电效应规律与方程、三个光电效应实验电路及Ek-ν/Uc-ν/I-U图像
%%%BLOCK id=005-01 ch=17 sec=热辐射·黑体辐射 kind=knowledge alt=none cont=none y=0.01-0.22
\textbf{选5} % 页面左上角的小标注

\noindent 热辐射 $\propto\left\{\begin{array}{@{}l@{}}\text{温度}\\ \text{材料}\\ \text{\unclear{表面}状况}\end{array}\right.$\quad （任一物体）辐射电磁波 % CHECK: “∝”手写形似“α”，按语义转为∝；括号内第三项“表面状况”被覆写，字迹不清

\noindent 黑体：吸收所有的入射电磁波而不反射的物体（小孔）\quad 不黑：太阳、窗户\quad 只辐射不反射

\noindent 对于$\langle$绝对黑体$\rangle$热辐射$\langle$只与$T$有关$\rangle$

\noindent $T\uparrow$\ 强度$\uparrow$\quad 峰值向短波移动

\noindent （红的发紫）

%FIG p005_a 0.40 0.122 0.585 0.197
\notefig[0.3]{figures/p005_a.png}
%%%END
%%%BLOCK id=005-02 ch=14 sec=电磁波谱 kind=summary alt=17 cont=none y=0.235-0.30
\noindent 无线电波\quad 微波\quad 红外线\,可见光\textsubscript{\,红\ 紫}\quad 紫外线\quad X射线\textsuperscript{(伦琴)}\quad $\gamma$射线\quad $\longrightarrow$\quad $f$大\ $\lambda$小 % 原稿：自左向右的长箭头画在整行文字上方，箭头末端写“f大”，其下写“λ小”；“(伦琴)”在“X射线”上方，“红 紫”在“可见光”下方
%%%END
%%%BLOCK id=005-03 ch=17 sec=黑体辐射·普朗克能量子 kind=knowledge alt=none cont=none y=0.28-0.315
\noindent 普朗克\ $\varepsilon=h\nu$（能量子）\quad $h=6.63\times10^{-34}\unit{J\cdot s}$\quad 不连续$=$量子化
%%%END
%%%BLOCK id=005-04 ch=17 sec=氢原子光谱·可见光谱线 kind=note alt=none cont=none y=0.255-0.335
\noindent $\langle$H光谱只有4种可见光\\
\hspace*{1.5em}$6\to2$\quad $5\to2$\\
\hspace*{1.5em}$4\to2$\quad $3\to2$ $\rangle$ % 原稿：两侧为一对大括号/尖括号
%%%END
%%%BLOCK id=005-05 ch=17 sec=光子能量·电子伏特 kind=knowledge alt=09 cont=none y=0.255-0.375
\noindent 电子 $1e$ 经 $U$ 后\\
\hspace*{2em}$W=-eU$\quad 自带负电

\noindent 电子伏特\quad $1e\cdot1\mathrm{V}=1.6\times10^{-19}\mathrm{C}\cdot1\mathrm{V}=1.6\times10^{-19}\mathrm{J}=1\mathrm{eV}$

\[
\underset{\text{光子能量}}{\text{可见光范围}}\quad \overset{\text{红外}}{1.6\unit{eV}}-\overset{\text{紫外}}{3.1\unit{eV}}
\] % 原稿：“光子能量”写在“可见光范围”下方；“红外”“紫外”两个小字分别写在“1.6eV”“3.1eV”上方(位置判断，可核对)
%%%END
%%%BLOCK id=005-06 ch=17 sec=光电效应·规律与方程 kind=knowledge alt=none cont=none y=0.35-0.60
\begin{itemize}
\item 赫兹$\langle$实验\unclear{发现}$\rangle$光电效应 % 原稿：“发现”两字与其它笔迹重叠，不清；该行与下一行前有大括号
\item 爱因斯坦：引入光子，解释光电效应（粒子性）\quad $\varepsilon=h\nu=h\dfrac{c}{\lambda}$
\end{itemize}

\noindent 光强/功率\quad $P=n\cdot h\nu$\quad （$n$为$\overset{\text{1s}}{\text{单位时间}}$$\overset{\text{光电子数}}{\text{光子数}}$） % 原稿：“单位时间”上方小注“1s”，“光子数”上方小注“光电子数”

\noindent 逸出功\quad 由金属决定\quad 同一金属 $W_0$ 相同

\noindent 光子（光照能量）$h\nu=h\dfrac{c}{\lambda}$\quad 仅由 $\nu$ 决定\quad $\neq$颜色

\noindent $\langle$最大初动能$\rangle$ $E_k=h\nu-W_0=eU_c$\quad 当 $h\nu=W_0$ 时\quad $\nu=\dfrac{W_0}{h}\to$ 无光电子/光电流

\noindent $\nu_0=\dfrac{W_0}{h}$ $\underset{\text{(极限频率)}}{\text{为截止频率}}$\quad 由金属决定\quad $\nu<\nu_0$ 时 光电流为0\quad 一个光子对应一个光电子

\noindent 光子 $\leftrightarrow$ 电子 一对一\quad
\begin{tabular}[t]{@{}lll@{}}
 & & 强蓝光\ 弱\unclear{蓝}光\\
个体 & 频率---决定初动能 & $\nu$相同\\
整体 & 光强---决定电子数/光子数 & $P$\textsubscript{光强}不同\\
\end{tabular} % 原稿：“个体”“整体”后接大括号；右侧“强蓝光 弱蓝光/ν相同/P(光强)不同”为小注

\noindent 发生条件\quad $h\nu>W_0$ 或 $\nu>\nu_0$
%%%END
%%%BLOCK id=005-07 ch=17 sec=光电效应·实验电路 kind=summary exp=1 alt=none cont=none y=0.595-0.64
\noindent \textbf{$\langle$三个电路$\rangle$}

\noindent 阻止\ （\unclear{光向}正极板）\ 光强不变，仅提高照射光频率，光子数小，电流$\downarrow$ % “光向”两字被覆写，不清；该行位于电路②③标题上方
%%%END
%%%BLOCK id=005-08 ch=17 sec=光电效应·实验电路 kind=diagram exp=1 alt=none cont=none y=0.635-0.745
\noindent \red{①} 光电效应发生电路

%FIG p005_b 0.125 0.6625 0.275 0.7395
\notefig[0.3]{figures/p005_b.png}
\noindent 瞬时性（不要时间） % CHECK: 首字手写形似“目急”，按语义转为“瞬”
%%%END
%%%BLOCK id=005-09 ch=17 sec=光电效应·规律与方程 kind=knowledge alt=none cont=none y=0.74-0.88
\noindent \textbf{发生要求}
\begin{itemize}
\item 大于极限频率
\item 小于 最大波长
\end{itemize}

\noindent 同一种光 $U_c$ 相同\\
同一种金属 $W_0$ 相同
%%%END
%%%BLOCK id=005-10 ch=17 sec=光电效应·实验电路 kind=diagram exp=1 alt=none cont=none y=0.635-0.90
\noindent \red{②} 遏止电压 $U_c$ 测定电路

%FIG p005_c 0.355 0.661 0.59 0.7645
\notefig[0.35]{figures/p005_c.png}
\noindent 负极接收\textsuperscript{极}$\to$遏止\\ % 原稿：“极”字小写在“收”上方，为插入
看到遏止电极压 $U_c$ % CHECK: 疑为“遏止电压”，“极”字原稿可见，照录
\[
\begin{array}{l}
-eU_c=0-E_k\\
\hookrightarrow\ eU_c=E_k=h\nu-W_0=h\dfrac{c}{\lambda}-W_0
\end{array}
\]
%%%END
%%%BLOCK id=005-11 ch=17 sec=光电效应·实验电路 kind=diagram exp=1 alt=none cont=none y=0.635-0.90
\noindent \red{③} \underline{饱和光电流测定电路}\ 光向负极板

%FIG p005_d 0.66 0.666 0.8235 0.764
\notefig[0.3]{figures/p005_d.png}
\noindent 原因：单位时间 $K$ 发出光电子数一定，达到某个电压后，光电子全部\unclear{被}阳极吸收，电流不再改变 % 原稿：“原因”二字写在图中阳极板右上角；“阳极”为行间小字插入；图中“hν”处引出一条弧线箭头指向“单位时间K”

\noindent 正极接收板$\to$促进\\
饱和光电流 只和 $n$ 有关\ （光子数）

\noindent 光强 $P=nh\nu$\\
$n$ 单位时间光子数\ （电子数）
%%%END
%%%BLOCK id=005-12 ch=17 sec=光电效应·图像 kind=diagram alt=none cont=none y=0.88-0.99
%FIG p005_e 0.015 0.8805 0.218 0.978
\notefig[0.3]{figures/p005_e.png}
\noindent $E_k=h\cdot\nu-W_0$ % 以上图中已含(含纵轴截距标注“-W_0”)

%FIG p005_f 0.218 0.8825 0.425 0.99
\notefig[0.3]{figures/p005_f.png}
\noindent $U_c=\dfrac{h}{e}\cdot\nu-\dfrac{W_0}{e}$ % 以上图中已含(含纵轴截距标注“-W_0/e”)

%FIG p005_g 0.425 0.89 0.592 0.962
\notefig[0.3]{figures/p005_g.png}
\noindent 高强光 左大频 % CHECK: 同第2页，首字形似“高”，语义疑为“同”
%%%END
%%%PAGEEND 5
